\writeSectionFrame{\presentationTitle}{Das Problem}
\begin{frame}{Onlineshop}
	\begin{center}
 		\includegraphics[scale=0.3]{\BILDERPFAD/problem1.png}
 	\end{center}
\end{frame}

\note{
	Ein Onlineshop verkauft verschiedene Waren. Natürlich sieht man
	relativ schnell, das man hier Vererbung einsetzen kann, um Gemeinsamkeiten
	herauszufaktorisieren.
}

\begin{frame}{Onlineshop}
	\begin{center}
 		\includegraphics[scale=0.25]{\BILDERPFAD/problem2.png}
 	\end{center}
\end{frame}

\note{
	
}

\begin{frame}{Onlineshop}
	\begin{center}
 		\includegraphics[scale=0.25]{\BILDERPFAD/problem3.png}
 	\end{center}
\end{frame}

\note{
	Nun soll die Artikelklasse noch um die Möglichkeit erweitert werden, dass die Artikel in
	verschiedene Formate, namentlich XML und pdf. Außerdem soll die Spiele-Klasse erweitert
	werden um eine Methode für die Onlineregistrierung. Das Problem ist, dass die Artikelklassen
	im Produktivbetrieb sind und da sie für den Betrieb des Shops nun einmal extrem wichtig sind,
	dürfen wir keine Änderungen daran vornehmen. Zudem ist jedem klar, dass es 
	wahrscheinlich nicht bei diesen beiden Export-Methoden bleiben wird. Wenn jetzt noch in HTML 
	exportiert werden soll, müssen wir die Artikelklassen noch einmal anfassen. 
}